\section{Conventions}\label{sec:conv}

\begin{enumerate}[(1),leftmargin=*]
\item An $h\times m$ grid has $h$ rows, $m$ columns, and point set $V$. A $2\times m$ grid
is a \emph{strip}. The point $(i,j)$ lies in column $i$ and row $j$, with $(1,1)$ at the
top left (Figure~\ref{fig:game}(a)).
\item A \emph{rectangle} $R$ is its set $C(R)$ of four corners, grid points spanning a
rectangle of positive area. It is \emph{axis-aligned} if its sides are parallel to the
grid lines, and \emph{tilted} otherwise.
\item A \emph{blackout} is a set $S\subseteq V$ with \emph{kept set} $K=V\setminus S$. The
\emph{presentation} of $R$ is $\pres(R)=C(R)\setminus S$.
\item Column $i$ has \emph{black rows} $B_i=\{j:(i,j)\in S\}$, and is \emph{empty}, a
\emph{singleton}, a \emph{double}, or a \emph{triple} when $|B_i|$ is $0$, $1$, $2$, or
$3$, and \emph{full} when $|B_i|=h$. The \emph{pair set} $F$ is the set of row pairs
contained in some $B_i$, and $q_h=\lfloor h^2/4\rfloor$ is \seqnum{A002620}$(h)$.
\end{enumerate}

\begin{definition}[valid blackout]\label{def:valid}
A blackout $S$ is \emph{valid} if $\pres(R_1)\ne\pres(R_2)$ whenever $R_1\ne R_2$. The
\emph{maximum blackout} $f(h,m)$ is the largest size of a valid blackout of the
$h\times m$ grid, and $a_h(m)$ counts the valid blackouts of that size, with their
coordinates.
\end{definition}

By symmetry $f(h,m)=f(m,h)$ and $a_h(m)=a_m(h)$. Width $0$ and height $1$ have no
rectangles, so $f(h,0)=0$, $a_h(0)=1$, $f(1,m)=m$, and $a_1(m)=1$. An empty presentation
is allowed, so a valid blackout hides at most one rectangle.

\begin{lemma}[difference sets]\label{lem:diff}
For sets $A$, $B$, and $S$, we have $A\setminus S=B\setminus S$ if and only if
$A\triangle B\subseteq S$.
\end{lemma}

\begin{proof}
Both conditions say that no point outside $S$ lies in exactly one of $A$ and $B$.
\end{proof}

\begin{corollary}[hitting sets]\label{cor:hitting}
By Lemma~\ref{lem:diff}, a blackout $S$ is valid if and only if $K$ meets $C(R_1)\triangle C(R_2)$ for all distinct
rectangles $R_1$ and $R_2$.
\end{corollary}

So subsets of valid blackouts are valid, and showing that no blackout of $k$ points is
valid proves $f(h,m)<k$.

\begin{lemma}[three corners suffice]\label{lem:threecorners}
Three corners of a rectangle determine it, so distinct rectangles with the same
presentation show at most two corners.
\end{lemma}

\begin{proof}
Three corners form a right triangle, whose only right angle is at the corner $q$ adjacent
to the other two, $p$ and $r$. So the fourth corner is $p+r-q$.
\end{proof}
